\documentclass[12pt,a4paper]{article}
\usepackage[utf8]{inputenc}
\usepackage[french]{babel}
\usepackage[T1]{fontenc}
\usepackage{geometry}
\geometry{margin=2.5cm}
\usepackage{array}
\usepackage{titlesec}
\usepackage{hyperref}
\usepackage{enumitem}
\usepackage{xcolor}

\titleformat{\section}{\large\bfseries}{\thesection}{1em}{}[\titlerule]

\title{\textbf{Business Plan : Obscurus Anti-Cheat Solutions}}
\author{Maël - Fondateur \& CTO}
\date{Mars 2026}

\begin{document}

\maketitle
\tableofcontents
\newpage

\section{Genèse et Vision du Projet}
Le projet Obscurus naît du constat que les solutions anti-triche actuelles sont soit trop intrusives (Kernel obligatoire), soit inefficaces face aux nouvelles menaces (DMA, IA). Obscurus propose une rupture technologique par l'hybridation : une protection adaptative qui évolue avec le cycle de vie du jeu vidéo.

\section{Organisation des Ressources Humaines}
L'entreprise est une SARL organisée pour maximiser la réactivité technique et commerciale.

\begin{itemize}[label=\textbullet]
    \item \textbf{Maël (PDG / CTO) :} 
    \begin{itemize}
        \item Développement du cœur système (C++/ASM) et du pilote Kernel.
        \item Architecture des micro-services Cloud (Go/Rust) pour l'analyse en temps réel.
        \item Vision technologique et roadmap R\&D.
    \end{itemize}
    \item \textbf{Anna (Responsable Cybersécurité / RSO) :} 
    \begin{itemize}
        \item Direction des missions "Red Team" (Reverse-engineering des binaires clients).
        \item Analyse de la menace et création de signatures de détection.
        \item Veille sur le "Dark Web" et les forums spécialisés en création de cheats.
    \end{itemize}
    \item \textbf{Joan (Responsable Commercial \& Communication) :} 
    \begin{itemize}
        \item Négociation des contrats B2B avec les studios et éditeurs.
        \item Stratégie de "Content Marketing" pour crédibiliser la solution auprès des joueurs.
        \item Gestion des relations publiques et salons (Gamescom, GDC).
    \end{itemize}
    \item \textbf{Camille (Responsable RH \& Comptabilité) :} 
    \begin{itemize}
        \item Gestion de la facturation récurrente (SaaS) et des bilans financiers.
        \item Aspects juridiques : conformité RGPD (protection des données joueurs) et EULA.
        \item Recrutement des futurs ingénieurs support et développeurs.
    \end{itemize}
\end{itemize}

\section{Architecture Technique et Offre Modulaire}
Obscurus ne vend pas un produit unique, mais une suite de services activables à la demande.

\subsection{Phase 1 : Obscurus Cloud \& IA}
\begin{itemize}
    \item \textbf{Moteur comportemental :} Analyse de la télémétrie (positions, angles de vue, inputs).
    \item \textbf{Module IA (Premium) :} Réseaux de neurones entraînés à détecter le "Smoothing" des Aimbots et les micro-corrections de trajectoire non-humaines.
    \item \textbf{Infrastructure :} Déploiement sur AWS/Azure avec une latence d'analyse < 50ms.
\end{itemize}

\subsection{Phase 2 : Obscurus Kernel Shield}
\begin{itemize}
    \item \textbf{Protection Ring 0 :} Pilote de périphérique bloquant l'accès à la mémoire du jeu par des processus tiers non autorisés.
    \item \textbf{Anti-Debugging :} Techniques avancées pour empêcher les tricheurs d'analyser le fonctionnement d'Obscurus.
\end{itemize}

\subsection{Prestation Audit Red Team}
\begin{itemize}
    \item Test d'intrusion applicatif (Pentest).
    \item Rapport détaillé des failles de logique de jeu (ex: duplication d'items, bypass de boutique).
    \item Recommandations de correction de code.
\end{itemize}

\section{Stratégie de Développement et Plan d'Action}
\begin{enumerate}
    \item \textbf{Mois 1-4 :} Développement du MVP Cloud et du Dashboard de gestion. Anna réalise les premiers audits rémunérés pour financer l'infrastructure.
    \item \textbf{Mois 5-8 :} Phase de bêta-test avec des studios partenaires. Entraînement de l'IA sur des données réelles. Joan lance la campagne de pré-commercialisation.
    \item \textbf{Mois 9-15 :} Lancement de la version Kernel et certification WHQL par Microsoft. Camille structure les contrats de support 24/7.
\end{enumerate}

\section{Analyse Financière Prévisionnelle (Résumé)}
\begin{itemize}
    \item \textbf{Coûts fixes :} Salaires, infrastructure Cloud, certifications de sécurité.
    \item \textbf{Revenus :} Abonnements SaaS (au nombre de joueurs actifs), Prestations d'audit (forfaitaire), Support premium.
\end{itemize}

\end{document}